\section{Cấu hình thí nghiệm}
\label{sec:setup}

\paragraph{Hạ tầng tính toán.}
Toàn bộ huấn luyện paper-faithful được thực hiện trên Kaggle Notebook GPU:
\begin{itemize}\itemsep0pt
    \item GPU: NVIDIA Tesla P100-PCIE-16GB (sm\_60, Pascal).
    \item Python 3.10, PyTorch 2.5.1+cu121 (downgrade từ default Kaggle 2.10 để hỗ trợ sm\_60).
    \item Mixed precision fp16 (AMP) cho phép batch 8 trong giới hạn 16GB VRAM.
\end{itemize}

\paragraph{Hyperparameters.}
Theo cấu hình paper sect.~5.2 (giữ nguyên với mọi biến thể để đảm bảo so sánh công bằng):
\begin{center}
\begin{tabular}{ll}
\toprule
Tổng epoch & 120 (Pha I: 40, Pha II: 80) \\
Optimizer  & 6 Adam riêng (Encoder, Decoder, BaseFuseLayer, DetailFuseLayer, AG-Base, AG-Detail) \\
Learning rate & $10^{-4}$, StepLR $\gamma = 0.5$ mỗi 20 epoch, floor $10^{-6}$ \\
Batch size & 8, AMP fp16 \\
Clip grad norm & 0.01 \\
$\alpha_1, \alpha_2, \alpha_3, \alpha_4$ & $1, 2, 5, 2$ (theo code release) \\
Seed & 42 \\
\bottomrule
\end{tabular}
\end{center}

\paragraph{Dữ liệu.}
\begin{itemize}\itemsep0pt
    \item \textbf{Huấn luyện}: 130 cặp Harvard medical theo split paper, patch $128 \times 128$ stride 64 (sau filter low-contrast còn $\sim 1100$ patches).
    \item \textbf{Đánh giá}: 72 cặp test (24 mỗi modality CT, PET, SPECT) từ \texttt{data/reference/}, dùng chung với các phương pháp SOTA trong Bảng~\ref{tab:sota-raw-metrics}.
\end{itemize}

\paragraph{Baseline so sánh.}
Để đảm bảo \emph{so sánh apples-to-apples}, đồ án không chỉ so với CDDFuse pretrained gốc mà còn \emph{huấn luyện lại} CDDFuse từ đầu với cùng quy trình paper-faithful trên dữ liệu của đồ án (gọi là \textbf{CDDFuse-Paper-MIF}). So sánh chính của đồ án là:
\begin{equation*}
\boxed{\textbf{CDDFuse-AG} \quad \text{vs} \quad \textbf{CDDFuse-Paper-MIF}}
\end{equation*}
Hai mô hình dùng cùng hạ tầng, cùng dữ liệu, cùng split, cùng hyperparameters $\to$ chênh lệch chỉ phản ánh đóng góp của Adaptive Gating + Saliency-guided Pixel.

\section{Kết quả chính: CDDFuse-AG vs CDDFuse-Paper-MIF}
\label{sec:results-main}

\subsection{Số liệu trung bình theo modality}
\label{ssec:per-modal-numbers}

Bảng~\ref{tab:cddfuse-ag-per-modal} trình bày giá trị trung bình của 7 chỉ số quan trọng nhất trên từng modality. Giá trị \textbf{in đậm} là giá trị tốt hơn giữa hai mô hình.

\begin{table}[h]
\centering
\small
\caption{Giá trị trung bình theo modality. $\uparrow$ cao tốt hơn, $\downarrow$ thấp tốt hơn. \textbf{In đậm} = tốt hơn.}
\label{tab:cddfuse-ag-per-modal}
\setlength{\tabcolsep}{5pt}
\begin{tabular}{l|cc|cc|cc}
\toprule
& \multicolumn{2}{c|}{\textbf{CT}} & \multicolumn{2}{c|}{\textbf{PET}} & \multicolumn{2}{c}{\textbf{SPECT}} \\
\textbf{Metric} & Baseline & CDDFuse-AG & Baseline & CDDFuse-AG & Baseline & CDDFuse-AG \\
\midrule
SSIM $\uparrow$   & 1.4434          & \textbf{1.4907} & \textbf{1.3235} & 1.3166          & 1.4031          & \textbf{1.4523} \\
PSNR $\uparrow$   & 55.32           & \textbf{55.35}  & \textbf{55.26}  & 55.17           & 57.65           & \textbf{57.79}  \\
QG   $\uparrow$   & 0.6183          & \textbf{0.6224} & 0.7638          & \textbf{0.7861} & 0.7503          & \textbf{0.7656} \\
QM   $\uparrow$   & 0.0046          & \textbf{0.0051} & 0.1230          & \textbf{0.1243} & 0.0788          & \textbf{0.0808} \\
NABF $\downarrow$ & 0.0323          & \textbf{0.0300} & 0.0153          & \textbf{0.0113} & \textbf{0.0194} & 0.0235          \\
QSF  $\uparrow$   & \textbf{0.0488} & 0.0307          & \textbf{0.0546} & 0.0506          & 0.0460          & \textbf{0.0495} \\
QMI  $\uparrow$   & 0.7830          & \textbf{0.7935} & \textbf{0.8020} & 0.8000          & 0.9049          & \textbf{0.9422} \\
\bottomrule
\end{tabular}
\end{table}

Bảng~\ref{tab:cddfuse-ag-pct} dưới đây thể hiện cùng số liệu dưới dạng \textbf{phần trăm cải thiện} so với baseline. Cell \textbf{in đậm} là cải thiện $> +3\%$ (đáng chú ý).

\begin{table}[h]
\centering
\small
\caption{Phần trăm cải thiện của CDDFuse-AG so với CDDFuse-Paper-MIF. \textbf{In đậm} = cải thiện $> +3\%$. Đỏ = giảm $> 3\%$ (chữ thường, không bold).}
\label{tab:cddfuse-ag-pct}
\begin{tabular}{lccc}
\toprule
\textbf{Metric} & \textbf{CT} & \textbf{PET} & \textbf{SPECT} \\
\midrule
SSIM   & \textbf{+3.3\%} & $-0.5\%$            & \textbf{+3.5\%} \\
PSNR   & $+0.1\%$        & $-0.2\%$            & $+0.2\%$ \\
QG     & $+0.7\%$        & $+2.9\%$            & $+2.0\%$ \\
QM     & \textbf{+9.5\%} & $+1.1\%$            & $+2.6\%$ \\
NABF (↓) & \textbf{+7.1\%} (đỡ artifact) & \textbf{+26.6\%} (đỡ artifact) & $-21.0\%$ (tệ hơn) \\
QSF    & $-37.1\%$       & $-7.3\%$            & \textbf{+7.5\%} \\
QMI    & $+1.4\%$        & $-0.3\%$            & \textbf{+4.1\%} \\
\bottomrule
\end{tabular}
\end{table}

\paragraph{Quan sát chính.}
\begin{itemize}\itemsep0pt
    \item \textbf{CT}: CDDFuse-AG \emph{thắng 5/7 metric chính}. Đáng chú ý SSIM $+3.3\%$, QM $+9.5\%$, NABF giảm $7.1\%$. Mất QSF ($-37\%$) là trade-off (nhường sharpness lấy SSIM/wavelet).
    \item \textbf{PET}: thắng 4/7. NABF cải thiện rất mạnh ($+26.6\%$), gradient tăng $+2.9\%$. Thua nhẹ SSIM và QSF.
    \item \textbf{SPECT}: thắng 6/7. SSIM $+3.5\%$, QSF $+7.5\%$, QMI $+4.1\%$ đều tăng. Mất NABF ($-21\%$) là chỗ duy nhất tệ hơn.
\end{itemize}

\subsection{Trung bình tổng quan trên toàn tập (Pooled)}
\label{ssec:pooled-results}

Khi gộp toàn bộ 72 cặp test (24 mỗi modality), giá trị trung bình \emph{toàn tập} của hai mô hình:

\begin{table}[h]
\centering
\small
\caption{Trung bình trên 72 cặp test gộp 3 modality. \textbf{In đậm} = tốt hơn.}
\label{tab:cddfuse-ag-pooled}
\begin{tabular}{lccc}
\toprule
\textbf{Metric} & \textbf{Baseline} & \textbf{CDDFuse-AG} & \textbf{Cải thiện (\%)} \\
\midrule
SSIM $\uparrow$   & 1.3900          & \textbf{1.4039} & $+1.00\%$  \\
PSNR $\uparrow$   & 56.0794         & \textbf{56.0872} & $+0.01\%$ (≈ tie) \\
QG   $\uparrow$   & 0.7108          & \textbf{0.7175} & $+0.94\%$ \\
QM   $\uparrow$   & \textbf{0.0688} & 0.0674          & $-2.05\%$ \\
NABF $\downarrow$ & 0.0224          & \textbf{0.0214} & \textbf{$+4.19\%$} (đỡ artifact) \\
QSF  $\uparrow$   & \textbf{0.0498} & 0.0447          & $-10.22\%$ \\
QMI  $\uparrow$   & 0.8300          & \textbf{0.8399} & $+1.20\%$ \\
MI   $\uparrow$   & \textbf{54.17}  & 51.71           & $-4.54\%$ \\
EN   $\uparrow$   & \textbf{5.0396} & 4.8923          & $-2.92\%$ \\
VAR  $\uparrow$   & \textbf{67.28}  & 65.93           & $-2.00\%$ \\
\bottomrule
\end{tabular}
\end{table}

\paragraph{Đánh giá tổng quan.}
CDDFuse-AG cải thiện \emph{nhất quán} ở các chỉ số quan trọng cho chẩn đoán y tế:
\begin{itemize}\itemsep0pt
    \item \textbf{NABF $\uparrow 4.2\%$}: ít ảo ảnh biên hơn --- mục tiêu chính của Adaptive Gating.
    \item \textbf{SSIM $\uparrow 1.0\%$}, \textbf{QMI $\uparrow 1.2\%$}, \textbf{QG $\uparrow 0.9\%$}: cấu trúc, mutual information, và gradient quality đều tốt hơn baseline.
\end{itemize}
Tuy nhiên có trade-off: MI, EN, VAR giảm $2\text{--}4\%$, QSF giảm $10\%$. Đây là biểu hiện của ``smoothing trade-off'' tự nhiên --- model học làm sạch biên (NABF) nhưng đánh đổi một phần thông tin nguồn (MI/VAR/EN).

\subsection{Kiểm định ý nghĩa thống kê}
\label{ssec:stats-results}

Áp dụng pipeline Wilcoxon + Cliff's $\delta$ + Holm--Bonferroni (xem Mục~\ref{sec:stats-bg}) trên 25 chỉ số chất lượng, kết quả \emph{số metric có ý nghĩa thống kê} (SIG):

\begin{table}[h]
\centering
\small
\caption{Số chỉ số có ý nghĩa thống kê sau Holm correction ($\alpha = 0.05$, $K = 22$). Per-modality test riêng từng modality ($n = 24$); Pooled gộp 3 modality ($n = 72$).}
\label{tab:cddfuse-ag-sig}
\begin{tabular}{lcccc}
\toprule
& \textbf{CT} & \textbf{PET} & \textbf{SPECT} & \textbf{Pooled} \\
\midrule
Số SIG / 22 metric & \textbf{10} & \textbf{10} & 2 & 2 \\
Verdict           & CONFIRM    & CONFIRM    & MIXED & MIXED \\
\bottomrule
\end{tabular}
\end{table}

\paragraph{Diễn giải.}
\begin{itemize}\itemsep0pt
    \item \textbf{CT, PET đạt CONFIRM\_IMPROVEMENT} ($\ge 5$ metric SIG): CDDFuse-AG \emph{cải thiện ý nghĩa} CDDFuse cho hai modality này.
    \item \textbf{SPECT chỉ 2 SIG}: hiệu quả không nhất quán. Có thể do baseline SSIM SPECT đã rất tốt (1.43) khó cải thiện hơn nữa, hoặc do đặc tính phân bố cường độ của SPECT.
    \item \textbf{Pooled 2 SIG}: gain mạnh ở CT/PET bị ``loãng'' bởi SPECT khi pool $\to$ tổng quan chỉ MIXED.
\end{itemize}

\section{Phân tích trade-off}
\label{sec:results-tradeoff}

Quan sát thấy CDDFuse-AG mạnh ở chỉ số \emph{cấu trúc + biên + saliency} (SSIM, QG, NABF, QMI) nhưng yếu hơn ở chỉ số \emph{thông tin tổng quát} (MI, EN, VAR, QSF). Điều này hợp lý với cơ chế:
\begin{itemize}\itemsep0pt
    \item Adaptive Gating \emph{lọc} các kênh feature theo trọng số mềm $\to$ giảm noise + artifact (NABF) nhưng cũng \emph{loại bỏ} một phần thông tin yếu $\to$ MI/VAR/EN giảm.
    \item Saliency-guided Pixel target ưu tiên vùng biên/texture (gradient cao) $\to$ giữ cấu trúc (SSIM) nhưng smooth vùng phẳng $\to$ QSF (sharpness toàn ảnh) giảm.
\end{itemize}
Đây không phải failure mà là \emph{lựa chọn thiết kế phù hợp với mục đích chẩn đoán}: bác sĩ quan tâm cấu trúc rõ và biên sạch hơn là độ phong phú entropy toàn ảnh.

\section{So sánh phụ với CDDFuse pretrained công bố}
\label{sec:results-vs-pretrained}

Để đặt CDDFuse-AG trong bối cảnh SOTA, đồ án cũng so sánh với checkpoint pretrained CDDFuse công bố (\texttt{CDDFuse\_MIF.pth}, không retrain). Đây không phải so sánh apples-to-apples (khác training data composition) nên chỉ là tham khảo.

\begin{table}[h]
\centering
\small
\caption{So với CDDFuse pretrained công bố --- tham khảo. \textbf{In đậm} = CDDFuse-AG tốt hơn pretrained.}
\label{tab:vs-pretrained}
\begin{tabular}{lccc}
\toprule
\textbf{Modal} & \textbf{Metric mạnh nhất} & \textbf{Cải thiện vs pretrained} & \textbf{Trade-off} \\
\midrule
PET   & \textbf{NABF} giảm \textbf{37\%}, \textbf{QM} tăng \textbf{68\%}, \textbf{QSF} tăng \textbf{54\%} & 3 metric large & MI, EN giảm \\
SPECT & \textbf{QSF} tăng \textbf{61\%}, \textbf{NABF} giảm \textbf{29\%}, \textbf{SSIM} tăng 1.4\% & 2 metric large & MI giảm \\
CT    & \textbf{SSIM} tăng 3.3\%, \textbf{QMI} tăng 2\% & nhỏ & QSF, MI giảm \\
\bottomrule
\end{tabular}
\end{table}

So với SOTA Bảng~\ref{tab:sota-raw-metrics}, CDDFuse-AG củng cố thêm vị trí top tier của CDDFuse, đặc biệt ở các chỉ số mà CDDFuse gốc còn yếu (NABF, QSF tùy modality).
